\section{ER-01 --- Identificador estruturado}
\subsection{Finalidade}
Reconhecer identificadores válidos da MiniLang.
\subsection{Alfabeto}
$\Sigma_{ID}=L\cup D\cup\{\_\}$, com $L$ para letras e $D$ para dígitos.
\subsection{Linguagem reconhecida}
A cadeia começa obrigatoriamente com uma letra; depois disso podem ocorrer caracteres alfanuméricos ou blocos iniciados por sublinhado, sendo que cada sublinhado deve ser seguido por pelo menos um caractere alfanumérico.
\subsection{ER formal}
$L(A\mid\_AA^*)^*$, com $A=L\cup D$.
\subsection{Sintaxe implementada em Python}
\begin{lstlisting}[language=Python]
r"[A-Za-z]([A-Za-z0-9]|_[A-Za-z0-9]+)*"
\end{lstlisting}
\subsection{Operadores utilizados}
A união $\mid$ permite escolher entre alternativas; a concatenação liga blocos em sequência; o fecho de Kleene $*$ permite zero ou mais repetições; e o fecho positivo $+$ indica uma ou mais ocorrências no ramo \texttt{\_[A-Za-z0-9]+}. Os parênteses agrupam subexpressões e as classes \texttt{[A-Za-z]} e \texttt{[A-Za-z0-9]} representam conjuntos finitos de símbolos.
\subsection{Equivalência formal e computacional}
\begin{center}
\small
\begin{tabular}{|p{0.43\textwidth}|p{0.43\textwidth}|}
\toprule
Formal & Python \\
\midrule
$L$ & \texttt{[A-Za-z]} \\
$A=L\cup D$ & \texttt{[A-Za-z0-9]} \\
$\_AA^*$ & \texttt{\_[A-Za-z0-9]+} \\
$+$ como $AA^*$ & \texttt{+} \\
\bottomrule
\end{tabular}
\end{center}
\subsection{AFN-epsilon}
Estado inicial: \texttt{q0}. Estado final: \texttt{q27}. A construção possui 28 estados, 35 transições, 27 transições $\varepsilon$ e 8 transições rotuladas.
Os diagramas JFLAP e Mermaid correspondentes estão em \texttt{docs/diagramas/}ER-01/. Ambos são derivados da mesma tabela canônica.
\subsection{Testes dirigidos}
\begin{center}
\small
\begin{tabular}{|p{0.34\textwidth}|c|p{0.43\textwidth}|}
\toprule
Cadeia & Esperado & Interpretação \\
\midrule
\texttt{a} & Aceita & Pertence à linguagem definida. \\
\texttt{nome} & Aceita & Pertence à linguagem definida. \\
\texttt{aluno1} & Aceita & Pertence à linguagem definida. \\
\texttt{valor\_total} & Aceita & Pertence à linguagem definida. \\
\texttt{mediaFinal2} & Aceita & Pertence à linguagem definida. \\
\texttt{a1\_b2} & Aceita & Pertence à linguagem definida. \\
$\varepsilon$ & Rejeita & Não pertence à linguagem definida. \\
\texttt{1aluno} & Rejeita & Não pertence à linguagem definida. \\
\texttt{\_nome} & Rejeita & Não pertence à linguagem definida. \\
\texttt{nome\_} & Rejeita & Não pertence à linguagem definida. \\
\texttt{nome\_\_aluno} & Rejeita & Não pertence à linguagem definida. \\
\texttt{nome-aluno} & Rejeita & Não pertence à linguagem definida. \\
\bottomrule
\end{tabular}
\end{center}
\subsection{Casos-limite}
Os casos de fronteira selecionados são: \texttt{a}, $\varepsilon$, \texttt{nome\_}, \texttt{nome\_\_aluno}.
\medskip
O propósito desses casos é explorar as menores cadeias válidas, limites superiores/inferiores e formas quase válidas que poderiam ser aceitas por uma regra permissiva demais.
\subsection{Observação sobre palavras reservadas}
Palavras como \texttt{int} e \texttt{float} possuem formato compatível com a ER de identificadores. A distinção final ocorre pela tabela de palavras reservadas do lexer, após o reconhecimento do lexema.

\section{ER-02 --- Inteiro sem zero à esquerda}
\subsection{Finalidade}
Reconhecer inteiros sem zero à esquerda, permitindo somente o literal `0` como zero isolado.
\subsection{Alfabeto}
$\Sigma_{INT}=D$, com $D=\{0,1,\ldots,9\}$.
\subsection{Linguagem reconhecida}
$L_{INT}=\{0\}\cup ND^*$, onde $N=\{1,\ldots,9\}$.
\subsection{ER formal}
$0\mid ND^*$.
\subsection{Sintaxe implementada em Python}
\begin{lstlisting}[language=Python]
r"(0|[1-9][0-9]*)"
\end{lstlisting}
\subsection{Operadores utilizados}
A união $\mid$ separa o caso do zero do caso iniciado por $N$; a concatenação forma a sequência dos símbolos; o fecho $*$ permite qualquer quantidade de dígitos após o primeiro não nulo; e os parênteses agrupam a alternativa no padrão Python.
\subsection{Equivalência formal e computacional}
\begin{center}
\small
\begin{tabular}{|p{0.43\textwidth}|p{0.43\textwidth}|}
\toprule
Formal & Python \\
\midrule
$0$ & \texttt{0} \\
$N$ & \texttt{[1-9]} \\
$D$ & \texttt{[0-9]} \\
$D^*$ & \texttt{[0-9]*} \\
\bottomrule
\end{tabular}
\end{center}
\subsection{AFN-epsilon}
Estado inicial: \texttt{q0}. Estado final: \texttt{q1}. A construção possui 10 estados, 12 transições, 9 transições $\varepsilon$ e 3 transições rotuladas.
Os diagramas JFLAP e Mermaid correspondentes estão em \texttt{docs/diagramas/}ER-02/. Ambos são derivados da mesma tabela canônica.
\subsection{Testes dirigidos}
\begin{center}
\small
\begin{tabular}{|p{0.34\textwidth}|c|p{0.43\textwidth}|}
\toprule
Cadeia & Esperado & Interpretação \\
\midrule
\texttt{0} & Aceita & Pertence à linguagem definida. \\
\texttt{1} & Aceita & Pertence à linguagem definida. \\
\texttt{7} & Aceita & Pertence à linguagem definida. \\
\texttt{9} & Aceita & Pertence à linguagem definida. \\
\texttt{10} & Aceita & Pertence à linguagem definida. \\
\texttt{20} & Aceita & Pertence à linguagem definida. \\
\texttt{2026} & Aceita & Pertence à linguagem definida. \\
\texttt{999} & Aceita & Pertence à linguagem definida. \\
$\varepsilon$ & Rejeita & Não pertence à linguagem definida. \\
\texttt{00} & Rejeita & Não pertence à linguagem definida. \\
\texttt{01} & Rejeita & Não pertence à linguagem definida. \\
\texttt{007} & Rejeita & Não pertence à linguagem definida. \\
\texttt{0123} & Rejeita & Não pertence à linguagem definida. \\
\texttt{-1} & Rejeita & Não pertence à linguagem definida. \\
\texttt{1.5} & Rejeita & Não pertence à linguagem definida. \\
\bottomrule
\end{tabular}
\end{center}
\subsection{Casos-limite}
Os casos de fronteira selecionados são: \texttt{0}, \texttt{00}, \texttt{01}, $\varepsilon$.
\medskip
O propósito desses casos é explorar as menores cadeias válidas, limites superiores/inferiores e formas quase válidas que poderiam ser aceitas por uma regra permissiva demais.

\section{ER-03 --- Decimal}
\subsection{Finalidade}
Reconhecer números decimais cuja parte inteira segue a regra de inteiro e cuja parte fracionária possui ao menos um dígito.
\subsection{Alfabeto}
$\Sigma_{DEC}=D\cup\{.\}$.
\subsection{Linguagem reconhecida}
A cadeia possui uma parte inteira válida, seguida por ponto e por uma ou mais ocorrências de dígitos.
\subsection{ER formal}
$(0\mid ND^*)PDD^*$, com $P=\{.\}$.
\subsection{Sintaxe implementada em Python}
\begin{lstlisting}[language=Python]
r"(0|[1-9][0-9]*)\.[0-9]+"
\end{lstlisting}
\subsection{Operadores utilizados}
A união $\mid$ define as alternativas da parte inteira; a concatenação coloca parte inteira, ponto e parte fracionária em sequência; o ponto é tratado como literal por meio de $P=\{.\}$ e do escape \texttt{\textbackslash.} no código; e o $+$ computacional exige pelo menos um dígito após o ponto.
\subsection{Equivalência formal e computacional}
\begin{center}
\small
\begin{tabular}{|p{0.43\textwidth}|p{0.43\textwidth}|}
\toprule
Formal & Python \\
\midrule
$N$ & \texttt{[1-9]} \\
$D$ & \texttt{[0-9]} \\
$P$ & \texttt{\textbackslash.} \\
$DD^*$ & \texttt{[0-9]+} \\
\bottomrule
\end{tabular}
\end{center}
\subsection{AFN-epsilon}
Estado inicial: \texttt{q0}. Estado final: \texttt{q17}. A construção possui 18 estados, 22 transições, 16 transições $\varepsilon$ e 6 transições rotuladas.
Os diagramas JFLAP e Mermaid correspondentes estão em \texttt{docs/diagramas/}ER-03/. Ambos são derivados da mesma tabela canônica.
\subsection{Testes dirigidos}
\begin{center}
\small
\begin{tabular}{|p{0.34\textwidth}|c|p{0.43\textwidth}|}
\toprule
Cadeia & Esperado & Interpretação \\
\midrule
\texttt{0.5} & Aceita & Pertence à linguagem definida. \\
\texttt{0.50} & Aceita & Pertence à linguagem definida. \\
\texttt{1.5} & Aceita & Pertence à linguagem definida. \\
\texttt{9.5} & Aceita & Pertence à linguagem definida. \\
\texttt{10.25} & Aceita & Pertence à linguagem definida. \\
\texttt{2026.75} & Aceita & Pertence à linguagem definida. \\
\texttt{00.5} & Rejeita & Não pertence à linguagem definida. \\
\texttt{01.5} & Rejeita & Não pertence à linguagem definida. \\
\texttt{.5} & Rejeita & Não pertence à linguagem definida. \\
\texttt{5.} & Rejeita & Não pertence à linguagem definida. \\
\texttt{1.2.3} & Rejeita & Não pertence à linguagem definida. \\
\texttt{-9.5} & Rejeita & Não pertence à linguagem definida. \\
\bottomrule
\end{tabular}
\end{center}
\subsection{Casos-limite}
Os casos de fronteira selecionados são: \texttt{0.5}, \texttt{5.}, \texttt{.5}, \texttt{00.5}.
\medskip
O propósito desses casos é explorar as menores cadeias válidas, limites superiores/inferiores e formas quase válidas que poderiam ser aceitas por uma regra permissiva demais.

\section{ER-04 --- Número científico}
\subsection{Finalidade}
Reconhecer números em notação científica com parte inteira válida, parte decimal opcional, expoente iniciado por `e` ou `E`, sinal opcional e expoente numérico obrigatório.
\subsection{Alfabeto}
$\Sigma_{SCI}=D\cup\{.,e,E,+,-\}$.
\subsection{Linguagem reconhecida}
A base é $0$ ou $ND^*$; a parte decimal é opcional; a letra $e$ ou $E$ é obrigatória; o sinal do expoente é opcional; o expoente possui um ou mais dígitos.
\subsection{ER formal}
$(0\mid ND^*)(PDD^*\mid\varepsilon)(e\mid E)(+\mid-\mid\varepsilon)DD^*$.
\subsection{Sintaxe implementada em Python}
\begin{lstlisting}[language=Python]
r"(0|[1-9][0-9]*)(\.[0-9]+)?[eE][+-]?[0-9]+"
\end{lstlisting}
\subsection{Operadores utilizados}
A união $\mid$ representa alternativas; a concatenação encadeia a mantissa, o expoente e seus componentes; a opcionalidade $?$ corresponde a $r\mid\varepsilon$ para a parte decimal; as classes \texttt{[eE]} e \texttt{[+-]} representam escolhas finitas; e \texttt{[0-9]+} exige um ou mais dígitos no expoente.
\subsection{Equivalência formal e computacional}
\begin{center}
\small
\begin{tabular}{|p{0.43\textwidth}|p{0.43\textwidth}|}
\toprule
Formal & Python \\
\midrule
$PDD^*\mid\varepsilon$ & \texttt{(\textbackslash.[0-9]+)?} \\
$e\mid E$ & \texttt{[eE]} \\
$+\mid-\mid\varepsilon$ & \texttt{[+-]?} \\
$DD^*$ & \texttt{[0-9]+} \\
\bottomrule
\end{tabular}
\end{center}
\subsection{AFN-epsilon}
Estado inicial: \texttt{q0}. Estado final: \texttt{q39}. A construção possui 42 estados, 52 transições, 40 transições $\varepsilon$ e 12 transições rotuladas.
Os diagramas JFLAP e Mermaid correspondentes estão em \texttt{docs/diagramas/}ER-04/. Ambos são derivados da mesma tabela canônica.
\subsection{Testes dirigidos}
\begin{center}
\small
\begin{tabular}{|p{0.34\textwidth}|c|p{0.43\textwidth}|}
\toprule
Cadeia & Esperado & Interpretação \\
\midrule
\texttt{0e0} & Aceita & Pertence à linguagem definida. \\
\texttt{1e10} & Aceita & Pertence à linguagem definida. \\
\texttt{9E9} & Aceita & Pertence à linguagem definida. \\
\texttt{1.5e3} & Aceita & Pertence à linguagem definida. \\
\texttt{12.25E-2} & Aceita & Pertence à linguagem definida. \\
\texttt{2026.75e+10} & Aceita & Pertence à linguagem definida. \\
\texttt{e10} & Rejeita & Não pertence à linguagem definida. \\
\texttt{1e} & Rejeita & Não pertence à linguagem definida. \\
\texttt{1.e3} & Rejeita & Não pertence à linguagem definida. \\
\texttt{.5e2} & Rejeita & Não pertence à linguagem definida. \\
\texttt{01e2} & Rejeita & Não pertence à linguagem definida. \\
\texttt{1.2e3.4} & Rejeita & Não pertence à linguagem definida. \\
\bottomrule
\end{tabular}
\end{center}
\subsection{Casos-limite}
Os casos de fronteira selecionados são: \texttt{0e0}, \texttt{1e}, \texttt{1.e3}, \texttt{01e2}.
\medskip
O propósito desses casos é explorar as menores cadeias válidas, limites superiores/inferiores e formas quase válidas que poderiam ser aceitas por uma regra permissiva demais.

\section{ER-05 --- Literal de string}
\subsection{Finalidade}
Reconhecer literais de string delimitados por aspas duplas, com conteúdo em um conjunto ASCII restrito.
\subsection{Alfabeto}
$C=L\cup D\cup\{\text{espaço},\texttt{\_},\texttt{.},\texttt{,},\texttt{!},\texttt{?},\texttt{:},\texttt{;},\texttt{+},\texttt{-},\texttt{*},\texttt{/},\texttt{=},\texttt{<},\texttt{>}\}$ e $\Sigma_{STR}=C\cup\{\text{aspas duplas}\}$.
\subsection{Linguagem reconhecida}
$L_{STR}=\{\text{aspas duplas }x\text{ aspas duplas}\mid x\in C^*\}$. A string vazia é válida.
\subsection{ER formal}
$\text{aspas duplas }C^*\text{ aspas duplas}$.
\subsection{Sintaxe implementada em Python}
\begin{lstlisting}[language=Python]
r'"[A-Za-z0-9 _.,!?;:+*/=<>-]*"'
\end{lstlisting}
\subsection{Operadores utilizados}
A classe finita do conteúdo representa os símbolos permitidos; o fecho de Kleene $*$ permite zero ou mais caracteres dentro das aspas; e as aspas duplas são delimitadores literais da cadeia. Não há mecanismo de escape nesta versão.
\subsection{Equivalência formal e computacional}
\begin{center}
\small
\begin{tabular}{|p{0.43\textwidth}|p{0.43\textwidth}|}
\toprule
Formal & Python \\
\midrule
$C^*$ & \texttt{[A-Za-z0-9 \_.,!?;:+*/=<>-]*} \\
aspas & \texttt{"..."} \\
$*$ permite $\varepsilon$ & \texttt{*} \\
\bottomrule
\end{tabular}
\end{center}
\subsection{AFN-epsilon}
Estado inicial: \texttt{q0}. Estado final: \texttt{q7}. A construção possui 8 estados, 9 transições, 6 transições $\varepsilon$ e 3 transições rotuladas.
Os diagramas JFLAP e Mermaid correspondentes estão em \texttt{docs/diagramas/}ER-05/. Ambos são derivados da mesma tabela canônica.
\subsection{Testes dirigidos}
\begin{center}
\small
\begin{tabular}{|p{0.34\textwidth}|c|p{0.43\textwidth}|}
\toprule
Cadeia & Esperado & Interpretação \\
\midrule
\texttt{\textquotedbl{}\textquotedbl{}} & Aceita & Pertence à linguagem definida. \\
\texttt{\textquotedbl{}Andrey\textquotedbl{}} & Aceita & Pertence à linguagem definida. \\
\texttt{\textquotedbl{}abc 123\textquotedbl{}} & Aceita & Pertence à linguagem definida. \\
\texttt{\textquotedbl{}a\_b.c\textquotedbl{}} & Aceita & Pertence à linguagem definida. \\
\texttt{\textquotedbl{}hello!\textquotedbl{}} & Aceita & Pertence à linguagem definida. \\
\texttt{\textquotedbl{}12:30\textquotedbl{}} & Aceita & Pertence à linguagem definida. \\
$\varepsilon$ & Rejeita & Não pertence à linguagem definida. \\
\texttt{\textquotedbl{}abc} & Rejeita & Não pertence à linguagem definida. \\
\texttt{abc\textquotedbl{}} & Rejeita & Não pertence à linguagem definida. \\
\texttt{\textquotedbl{}a\textquotedbl{}b\textquotedbl{}} & Rejeita & Não pertence à linguagem definida. \\
\texttt{\textquotedbl{}a\textbackslash{}b\textquotedbl{}} & Rejeita & Não pertence à linguagem definida. \\
\texttt{\textquotedbl{}linha} + quebra de linha & Rejeita & Não pertence à linguagem definida. \\
\bottomrule
\end{tabular}
\end{center}
\subsection{Casos-limite}
Os casos de fronteira selecionados são: \texttt{\textquotedbl{}\textquotedbl{}}, $\varepsilon$, \texttt{\textquotedbl{}abc}, \texttt{\textquotedbl{}a\textbackslash{}b\textquotedbl{}}.
\medskip
O propósito desses casos é explorar as menores cadeias válidas, limites superiores/inferiores e formas quase válidas que poderiam ser aceitas por uma regra permissiva demais.
\subsection{Limitação específica}
A primeira versão da MiniLang não possui mecanismo de escape. Aspas, barra invertida, tabulação e quebra de linha não pertencem ao conjunto de conteúdo \(C\).

\section{ER-06 --- Literal de horário}
\subsection{Finalidade}
Reconhecer literais de horário no formato HH:MM, com horas entre 00 e 23 e minutos entre 00 e 59.
\subsection{Alfabeto}
$\Sigma_{TIME}=D\cup\{:\}$.
\subsection{Linguagem reconhecida}
A hora pertence a 00--19 ou 20--23; os minutos pertencem a 00--59.
\subsection{ER formal}
$((0\mid1)D\mid2(0\mid1\mid2\mid3))Q(0\mid1\mid2\mid3\mid4\mid5)D$, com $Q=\{:\}$.
\subsection{Sintaxe implementada em Python}
\begin{lstlisting}[language=Python]
r"([01][0-9]|2[0-3]):[0-5][0-9]"
\end{lstlisting}
\subsection{Operadores utilizados}
A união $\mid$ define as alternativas para as horas e minutos; a concatenação organiza os dígitos e o separador; as classes \texttt{[01]}, \texttt{[0-3]} e \texttt{[0-5]} representam os conjuntos finitos correspondentes.
\subsection{Equivalência formal e computacional}
\begin{center}
\small
\begin{tabular}{|p{0.43\textwidth}|p{0.43\textwidth}|}
\toprule
Formal & Python \\
\midrule
$(0\mid1)$ & \texttt{[01]} \\
$D$ & \texttt{[0-9]} \\
$2(0\mid1\mid2\mid3)$ & \texttt{2[0-3]} \\
$Q$ & \texttt{:} \\
$(0\mid1\mid2\mid3\mid4\mid5)$ & \texttt{[0-5]} \\
\bottomrule
\end{tabular}
\end{center}
\subsection{AFN-epsilon}
Estado inicial: \texttt{q0}. Estado final: \texttt{q15}. A construção possui 16 estados, 16 transições, 9 transições $\varepsilon$ e 7 transições rotuladas.
Os diagramas JFLAP e Mermaid correspondentes estão em \texttt{docs/diagramas/}ER-06/. Ambos são derivados da mesma tabela canônica.
\subsection{Testes dirigidos}
\begin{center}
\small
\begin{tabular}{|p{0.34\textwidth}|c|p{0.43\textwidth}|}
\toprule
Cadeia & Esperado & Interpretação \\
\midrule
\texttt{00:00} & Aceita & Pertence à linguagem definida. \\
\texttt{09:05} & Aceita & Pertence à linguagem definida. \\
\texttt{12:30} & Aceita & Pertence à linguagem definida. \\
\texttt{14:59} & Aceita & Pertence à linguagem definida. \\
\texttt{23:00} & Aceita & Pertence à linguagem definida. \\
\texttt{23:59} & Aceita & Pertence à linguagem definida. \\
\texttt{24:00} & Rejeita & Não pertence à linguagem definida. \\
\texttt{12:60} & Rejeita & Não pertence à linguagem definida. \\
\texttt{2:30} & Rejeita & Não pertence à linguagem definida. \\
\texttt{14:5} & Rejeita & Não pertence à linguagem definida. \\
\texttt{99:99} & Rejeita & Não pertence à linguagem definida. \\
\texttt{14:30:00} & Rejeita & Não pertence à linguagem definida. \\
\bottomrule
\end{tabular}
\end{center}
\subsection{Casos-limite}
Os casos de fronteira selecionados são: \texttt{00:00}, \texttt{23:59}, \texttt{24:00}, \texttt{12:60}.
\medskip
O propósito desses casos é explorar as menores cadeias válidas, limites superiores/inferiores e formas quase válidas que poderiam ser aceitas por uma regra permissiva demais.
\subsection{Observação de correção}
Na faixa de horas 20--23, o segundo dígito pertence ao conjunto \(\{0,1,2,3\}\).

Nos arquivos JFLAP e Mermaid, essa transição aparece como \texttt{[0-3]}. 
